%! Setting Up --- Learning from Data — Unit 8, Lesson 8.0 — Exit Ticket (Answer Key)
\documentclass[10pt]{article}
\usepackage{linalg-article}
\usepackage{linalg-key}   % pulls in -boxes; do NOT also load -boxes
\newcommand{\TallMath}[1]{$\displaystyle #1\rule[-1.4em]{0pt}{3.2em}$}
\newcommand{\vv}{\mathbf{v}}
\newcommand{\relu}{\mathrm{ReLU}}

\begin{document}

\pageheader{Unit 8, Lesson 8.0}{Exit Ticket --- Answer Key}
\namedateperiod

\vspace{0.15in}

No notes. Show your reasoning.

\begin{enumerate}[leftmargin=1.8em, itemsep=16pt]
  \item \textbf{One layer.} For $A=\begin{bmatrix} 1 & -2 \\ 1 & 1 \end{bmatrix}$ and
        $\vv=\begin{bmatrix} 3 \\ 2 \end{bmatrix}$, compute the linear step $A\vv$, then apply the ReLU to
        each entry to get $\relu(A\vv)$.
        \ansline{$A\vv=\begin{bsmallmatrix} 3-4\\ 3+2 \end{bsmallmatrix}=\begin{bsmallmatrix} -1\\ 5 \end{bsmallmatrix}$, so $\relu(A\vv)=\begin{bsmallmatrix} 0\\ 5 \end{bsmallmatrix}$ (the $-1$ becomes $0$; the $5$ stays).}
  \item \textbf{Measure the error.} A network predicts $2,5,3$ while the true targets are $3,5,1$. Compute
        the loss (the sum of squared errors).
        \ansline{$(2-3)^2+(5-5)^2+(3-1)^2=1+0+4=5$.}
  \item \textbf{Explain.} In one sentence, why does a neural network need the nonlinear ReLU step ---
        what goes wrong if every step is just a matrix multiplication?
        \ansline{Without the ReLU, stacking linear steps just multiplies the matrices into one matrix --- a single straight-line function --- so the network could never bend to fit curved data.}
\end{enumerate}

\begin{teachernote}
Sort responses into ``got it / arithmetic slip / concept slip.'' Item~1 checks one layer: the linear
step $A\vv=\begin{bsmallmatrix} -1\\5 \end{bsmallmatrix}$ then ReLU \emph{entry-by-entry} to
$\begin{bsmallmatrix} 0\\5 \end{bsmallmatrix}$. Item~2 is the loss (sum of squared errors $=5$) --- watch
for students who forget to square or who drop the sign before squaring. Item~3 is the target idea: the
ReLU's nonlinearity is what keeps stacked layers from collapsing into a single matrix. If many stumble on
item~3, open Lesson~8.1 with a 30-second ``one ReLU $=$ one fold; folds let a function bend'' recap.
\end{teachernote}

\end{document}
